\chapter{Related Work: Comparative Studies and the Advantage Debate}
\label{ch:related}

The empirical part of this thesis enters a conversation with four strands of
literature: systematic benchmarks of quantum models, the ``power of data''
argument, classical surrogates and dequantization results, and the small
island of provable separations. This chapter reviews each and closes by
positioning the present work.

\section{Systematic benchmark studies}
\label{sec:related-benchmarks}

For most of QML's short history, empirical claims rested on one-dataset,
one-baseline demonstrations. \citet{bowles2024better} changed the standard:
benchmarking twelve popular quantum models (and classical controls) across
structured dataset families with systematic hyperparameter search, they found
--- almost uniformly --- that well-tuned classical models match or beat the
quantum contenders on standard classical-data tasks, that out-of-the-box
quantum models are highly sensitive to hyperparameters, and, strikingly, that
\emph{removing entanglement} from several quantum architectures often does not
degrade accuracy, suggesting the quantum ingredient frequently is not what
carries the performance. Their methodological catalogue of the field's
recurring sins --- tuned-quantum versus default-classical comparisons, absent
strong baselines (tree ensembles in particular), inconsistent splits, and
regime mixing --- reads as a checklist of what this thesis's fairness protocol
(Chapter~\ref{ch:methodology}) is designed to preclude. In the kernel corner
specifically, \citet{schnabel2025scrutiny} subjected fidelity- and
projected-kernel pipelines to the same discipline and reached similarly
sobering conclusions, while emphasizing bandwidth and preprocessing as
first-order factors --- directly informing the E2 design here. The
replication-style component of this thesis (entanglement on/off ablations, the
default-versus-tuned deltas of E1) is a deliberate test of whether these
findings transfer to a different dataset suite and protocol.

\section{The power of data}
\label{sec:related-powerofdata}

\citet{huang2021power} reframed the advantage question: many quantum
computations become classically learnable once a classical model has access
to \emph{data} generated by the process --- training data can substitute for
quantum computation. They provide a geometric, dataset-specific test (based on
the model-independent \emph{geometric difference} between kernel matrices)
that upper-bounds how much any quantum kernel can outperform a classical one
on a given dataset, and introduce the projected quantum kernel to engineer
tasks where the gap is provably large. Two lessons shape this thesis. First,
dataset-conditional diagnostics (alignment, geometric difference) belong in
the pipeline \emph{before} performance claims --- E2 computes them. Second,
the datasets on which quantum kernels could plausibly excel are special by
construction; finding an advantage on generic UCI-style data would contradict
theory, so the honest expectation for RQ1 is parity at best, and the
scientifically interesting content lies in \emph{where and why} deviations
occur (the sign-parity task being this thesis's engineered candidate).

\section{Classical surrogates and dequantization}
\label{sec:related-surrogates}

A complementary deflationary line shows that trained quantum models often
admit efficient \emph{classical surrogates}. \citet{schreiber2023surrogates}
exploit exactly the Fourier representation \eqref{eq:fourier}: since a
re-uploading model is a trigonometric polynomial with an
encoding-determined frequency set, one can fit that function class directly,
classically, from evaluations --- matching the quantum model's behavior
without a quantum device. More broadly, \citet{cerezo2023simulability} argue
that the very structures that provably avoid barren plateaus (shallow local
circuits, QCNNs, symmetry-restricted ansätze) confine the dynamics to
polynomially large subspaces that are classically simulable --- crystallizing
a dilemma: at NISQ scale, \emph{trainable} tends to imply \emph{simulable}.
For a thesis running at $\le 20$ qubits this is not a threat but a framing
device (Section~\ref{sec:intro-framing}): the object of study is the behavior
of these models, with simulability guaranteed either way.

\section{Provable separations and quantum data}
\label{sec:related-separations}

The existence proofs live elsewhere. \citet{liu2021rigorous} construct a
classification task from the discrete-logarithm problem for which a fidelity
quantum kernel learns efficiently while any efficient classical learner would
break a standard cryptographic assumption --- a rigorous, end-to-end
separation, on a deliberately cryptographic (not naturally occurring) task;
it defines the stretch experiment E10. On genuinely \emph{quantum} data, the
case is stronger still: \citet{huang2022advantage} demonstrate, on real
hardware, exponential advantages in learning from quantum-enhanced
experiments (QQ in the taxonomy). Both results delimit rather than contradict
this thesis: they identify where advantage is theoretically at home ---
structured algebraic tasks and quantum data --- and thereby explain why a
matched-conditions comparison on classical benchmarks (this work) should be
read as mapping the terrain of \emph{competitiveness}, not searching for
advantage where theory says none should live.

\section{Positioning of this thesis}
\label{sec:related-positioning}

Relative to this literature, the present work makes four commitments that,
combined, are not standard even in the strongest prior benchmarks:
(i)~\emph{one} stored split system --- 5-fold stratified $\times$ five seeds,
generated once, shared by every model of every family, forever; (ii)~a
\emph{matched and audited} tuning budget (one 50-trial study per
model--dataset pair under an identical procedure, with the full trial log
published); (iii)~\emph{parameter matching} for every NN-style comparison and
\emph{regime separation} (exact/shots/noisy/hardware never pooled) for every
reported number; and (iv)~a dataset suite deliberately spanning saturated to
hard, including an engineered parity-structured task that instantiates the
theory-motivated best case for entangling feature maps. The intended
contribution is thus less a verdict than a reproducible measurement
instrument --- in the spirit of \citet{bowles2024better}, extended with the
budget-matching and regime discipline their audit calls for.
